\section{Installation and Building}\label{installation-and-building}

The OpenUxAS build and dependency workflow changes as supported operating
systems and third-party dependencies evolve. The repository root
\texttt{README.md} is the canonical source for current platform requirements,
installation steps, and build commands.

Do not use the legacy Meson/Ninja installation instructions that were
previously reproduced in this user manual. OpenUxAS now provides the
\texttt{anod} command to obtain and build the required dependencies and
OpenUxAS itself.

\subsection{Quick start}\label{quick-start}

\begin{enumerate}
\item
  Check the repository root \texttt{README.md} for the currently supported
  operating systems and prerequisite tools.
\item
  Clone OpenUxAS and enter the repository:

\begin{verbatim}
git clone https://github.com/afrl-rq/OpenUxAS
cd OpenUxAS
\end{verbatim}

\item
  Fetch dependencies and build OpenUxAS:

\begin{verbatim}
./anod build uxas
\end{verbatim}

\item
  If OpenAMASE simulation support is required, build it with:

\begin{verbatim}
./anod build amase
\end{verbatim}

\item
  Run an example, for example:

\begin{verbatim}
./run-example 02_Example_WaterwaySearch
\end{verbatim}
\end{enumerate}

\subsection{Developing OpenUxAS}\label{developing-openuxas}

Run the \texttt{anod} build first so that the repository-local third-party
dependencies are available. After that initial build, the repository Makefile
can be used for incremental OpenUxAS builds:

\begin{verbatim}
make -j all
\end{verbatim}

If development also requires changes to LmcpGen or OpenAMASE, use the
development setup commands documented in the root \texttt{README.md}, for
example:

\begin{verbatim}
./anod devel-setup lmcp
./anod devel-setup amase
\end{verbatim}

\subsection{Running tests}\label{running-tests}

After OpenUxAS has been built, the C++ test suite can be run from
\texttt{tests/cpp}:

\begin{verbatim}
cd tests/cpp
./run-tests
\end{verbatim}

See \texttt{tests/cpp/README.md} for test-development guidance.

\subsection{Keeping installation guidance current}\label{keeping-installation-guidance-current}

The root \texttt{README.md} should be updated whenever supported platforms,
prerequisite tools, dependency management, or build commands change. This
user-manual page intentionally points to that canonical workflow rather than
duplicating platform-specific package lists that can become stale
independently.
